\chapter{Методологія дослідження}

\section{Джерела даних}

Дані отримано з публічного порталу Національного агентства із забезпечення
якості вищої освіти (НАЗЯВО): \url{https://public.naqa.gov.ua}.

Для кожної акредитаційної справи зі спеціальності 022 «Дизайн» було зібрано:
\begin{itemize}
    \item Форму СЕ з 16 вкладками (загальні відомості, критерії акредитації)
    \item Таблицю 1 --- перелік освітніх компонентів
    \item Таблицю 2 --- відомості про викладачів
    \item Силабуси навчальних дисциплін (PDF-файли)
\end{itemize}

\section{Збір даних (веб-скрейпінг)}

Для автоматизованого збору даних розроблено програмний інструмент на Python
з використанням бібліотеки Playwright для браузерної автоматизації.

\begin{lstlisting}[language=Python,caption={Архітектура скрейпера}]
# Основні модули:
# - navigator.py: навігація по порталу НАЗЯВО
# - extractor.py: витягування даних з форми СЕ
# - downloader.py: завантаження PDF-файлів
# - checkpoint.py: збереження прогресу
\end{lstlisting}

Скрейпер реалізує три фази:
\begin{enumerate}
    \item Збір URL-адрес акредитаційних справ
    \item Витягування даних з кожної справи
    \item Завантаження PDF-файлів (силабуси, звіти)
\end{enumerate}

\section{Обробка тексту}

Для обробки українськомовного тексту використано:
\begin{itemize}
    \item spaCy з українською моделлю (uk\_core\_news\_lg)
    \item pymorphy2 для лематизації
    \item Власний словник стоп-слів для академічних текстів
\end{itemize}

\section{Тематичне моделювання}

Застосовано BERTopic з мультимовною моделлю ембедингів
(paraphrase-multilingual-mpnet-base-v2).

Параметри моделювання:
\begin{itemize}
    \item min\_topic\_size = 10
    \item Автоматичне визначення кількості тем
    \item Кластеризація методом HDBSCAN
\end{itemize}

\section{Мережевий аналіз}

Побудовано граф спільної появи курсів:
\begin{itemize}
    \item Вершини: унікальні назви курсів (нормалізовані)
    \item Ребра: спільна присутність у навчальних планах
    \item Вага ребра: кількість програм зі спільними курсами
\end{itemize}

Для виявлення спільнот використано алгоритм Лувена.

\section{Аналіз прогалин у навичках}

Розроблено фреймворк AI та цифрових навичок для дизайнерів, що включає:
\begin{itemize}
    \item Базові AI-навички (основи ML, генеративний AI)
    \item 2D/3D цифрові інструменти (Adobe, Blender)
    \item UX/UI дизайн (Figma, прототипування)
    \item Візуалізація даних
    \item VR/AR дизайн
\end{itemize}

Покриття кожної навички оцінювалось за наявністю ключових слів у назвах
курсів та силабусах.
